\documentclass[11pt,a4paper]{article}

% Packages
\usepackage[utf8]{inputenc}
\usepackage[T1]{fontenc}
\usepackage{amsmath}
\usepackage{amsfonts}
\usepackage{amssymb}
\usepackage{geometry}
\usepackage{enumitem}
\usepackage{booktabs}
\usepackage{fancyhdr}
\usepackage{titlesec}

% Page setup
\geometry{margin=1in}
\setlength{\headheight}{13.6pt}
\pagestyle{fancy}
\fancyhf{}
\rhead{\thepage}
\lhead{Transformer Accounting}

% Title formatting
\titleformat{\section}{\large\bfseries}{\thesection}{1em}{}
\titleformat{\subsection}{\normalsize\bfseries}{\thesubsection}{1em}{}

% Document
\begin{document}

\title{Transformer Accounting: Simplifying Parameter and Compute Analysis}
\author{}
\date{\today}
\maketitle

\section{Introduction}

There are some insights that simplify thinking about the transformer quite a bit.

\section{Parameter Conventions}

The following conventions are used throughout this analysis:

\begin{align}
d_{ff} &= 4 \cdot d_{model} \\
d_q = d_k = d_v &= \frac{d_{model}}{n_{heads}}
\end{align}

We consider a non-gated MLP architecture.

\section{Parameter Structure}

\subsection{Matrix Organization}

There are 6 matrices with parameters per layer. They all have similar shapes:

\begin{itemize}
    \item \textbf{Two for MLP:}
    \begin{itemize}
        \item First matrix: $d_{model} \rightarrow d_{ff}$
        \item Second matrix: $d_{ff} \rightarrow d_{model}$
    \end{itemize}
    
    \item \textbf{Four for attention:}
    \begin{itemize}
        \item All have shape $d_{model} \rightarrow d_{model}$
    \end{itemize}
\end{itemize}

\subsection{Total Parameters}

The total number of parameters per layer is:

\begin{equation}
\text{Parameters} = 2 \cdot (d_{model} \cdot d_{ff}) + 4 \cdot d_{model}^2
\end{equation}

For the common case where $d_{ff} = 4 \cdot d_{model}$:

\begin{equation}
\text{Parameters} = 12 \cdot d_{model}^2
\end{equation}

In a sufficiently large model, this formula will get you within 1\% of non-embedding parameters.

\section{Computational Analysis}

\subsection{Matrix Multiplication Pattern}

Each parameter matrix is multiplied by each token embedding on each pass. The usage pattern is remarkably simple: all 6 matrices multiply the same inputs.

The input to a model has shape $[\text{batch}, \text{sequence}, d_{model}]$. It consists of $d_{model}$ vectors, where each vector is multiplied by each parameter matrix.

\subsection{Compute Requirements}

The total compute for multiplying with all matrices is:

\begin{equation}
\text{Compute} = 24 \cdot d_{model}^3 \cdot \text{seq} \cdot \text{batch}
\end{equation}

The compute per token is proportional to $2 \times \text{num\_params}$. Arithmetic intensity is directly linear to the number of tokens processed.

\section{Attention Mechanism}

There is additional compute required to handle interactions across the sequence, known as the attention mechanism.

Key properties of attention:
\begin{itemize}
    \item Does \textbf{not} add any parameters
    \item Compute is linear in $d_{model}$
    \item Compute is quadratic in sequence length
\end{itemize}

\section{Parameter and Compute Summary}

\begin{table}[h]
\centering
\begin{tabular}{|l|c|c|}
\hline
\textbf{Operation} & \textbf{Parameters} & \textbf{FLOPs per Token} \\
\hline
MLP Layer 1 & $d_{model} \times d_{ff}$ & $2d_{model} \times d_{ff}$ \\
\hline
MLP Layer 2 & $d_{ff} \times d_{model}$ & $2d_{ff} \times d_{model}$ \\
\hline
Attention: Q & $d_{model} \times d_{model}$ & $2d_{model}^2$ \\
\hline
Attention: K & $d_{model} \times d_{model}$ & $2d_{model}^2$ \\
\hline
Attention: V & $d_{model} \times d_{model}$ & $2d_{model}^2$ \\
\hline
Attention: Output & $d_{model} \times d_{model}$ & $2d_{model}^2$ \\
\hline
Attention: Scores & — & $2d_{model} \times \text{seq}^2$ \\
\hline
\textbf{Total per Layer} & $\mathbf{2d_{model} \times d_{ff} + 4d_{model}^2}$ & $\mathbf{2d_{model}(d_{ff} + 2d_{model}) + 2d_{model} \times \text{seq}^2}$ \\
\hline
\textbf{Common Case} & $\mathbf{12d_{model}^2}$ & $\mathbf{24d_{model}^2 + 2d_{model} \times \text{seq}^2}$ \\
\textbf{($d_{ff} = 4d_{model}$)} & & \\
\hline
\end{tabular}
\caption{Parameter count and compute requirements per transformer layer}
\end{table}

\end{document}
